\chapter{2020年新高考II卷}

\section{单选题}

\begin{problem}
设集合 \(A=\{2,3,5,7\}\)，\(B=\{1,2,3,5,8\}\)，则 \(A\cap B=\)（\quad）

\choices
    {\(\{1,3,5,7\}\)}
    {\(\{2,3\}\)}
    {\(\{2,3,5\}\)}
    {\(\{1,2,3,5,7,8\}\)}

\end{problem}

\begin{answer}
C
\end{answer}

\begin{solution}
由集合交集的定义可知
\(A\cap B=\{2,3,5\}\).故选 C.
\end{solution}

\begin{problem}
\((1+2\i)(2+\i)=\)（\quad）

\choices
    {\(4+5\i\)}
    {\(5\i\)}
    {\(-5\i\)}
    {\(2+3\i\)}

\end{problem}

\begin{answer}
B
\end{answer}

\begin{solution}
直接计算得
\((1+2\i)(2+\i)=2+\i+4\i+2\i^2=5\i\).故选 B.
\end{solution}

\begin{problem}
在 \(\triangle ABC\) 中，\(D\) 是 \(AB\) 边上的中点，则 \(\overrightarrow{CB}=\)（\quad）

\choices
    {\(2\overrightarrow{CD}+\overrightarrow{CA}\)}
    {\(\overrightarrow{CD}-2\overrightarrow{CA}\)}
    {\(2\overrightarrow{CD}-\overrightarrow{CA}\)}
    {\(\overrightarrow{CD}+2\overrightarrow{CA}\)}

\end{problem}

\begin{answer}
C
\end{answer}

\begin{solution}
因为 \(D\) 是 \(AB\) 边上的中点，所以
\(\overrightarrow{AB}=2\overrightarrow{AD}\).又因为
\(\overrightarrow{CB}=\overrightarrow{CA}+\overrightarrow{AB}\)，所以
\(\overrightarrow{CB}=\overrightarrow{CA}+2\overrightarrow{AD}\).而
\(\overrightarrow{CD}=\overrightarrow{CA}+\overrightarrow{AD}\)，故
\(\overrightarrow{AD}=\overrightarrow{CD}-\overrightarrow{CA}\).代入可得
\(\overrightarrow{CB}=\overrightarrow{CA}+2\bigl(\overrightarrow{CD}-\overrightarrow{CA}\bigr)=2\overrightarrow{CD}-\overrightarrow{CA}\).故选 C.
\end{solution}

\begin{problem}
日晷是中国古代用来测定时间的仪器，利用与晷面垂直的晷针投射到晷面的影子来测定时间. 把地球看成一个球（球心记为 \(O\)），地球上一点 \(A\) 的纬度是指 \(OA\) 与地球赤道所在平面所成角，点 \(A\) 处的水平面是指过点 \(A\) 且与 \(OA\) 垂直的平面. 在点 \(A\) 处放置一个日晷，若晷面与赤道所在平面平行，点 \(A\) 处的纬度为北纬 \(40^\circ\)，则晷针与点 \(A\) 处的水平面所成角为（\quad）

\begin{center}
\bitmapfigure[max width=0.42\linewidth,max totalheight=4.8cm]{img/2020/national_paper_2/q4_fig1.png}
\end{center}

\choices
    {\(20^\circ\)}
    {\(40^\circ\)}
    {\(50^\circ\)}
    {\(90^\circ\)}

\end{problem}

\begin{answer}
B
\end{answer}

\begin{solution}
画出过球心和晷针所确定的平面截地球和晷面的截面图，其中 \(CD\) 是赤道所在平面的截线，\(l\) 是点 \(A\) 处的水平面的截线，依题意可知 \(OA\perp l\)；\(AB\) 是晷针所在直线，\(m\) 是晷面的截线. 依题意，晷面和赤道平面平行，晷针与晷面垂直，根据平面平行的性质定理可得 \(m//CD\)，根据线面垂直的定义可得 \(AB\perp m\).

由于 \(\angle AOC=40^\circ\)，\(m//CD\)，所以
\(\angle OAG=\angle AOC=40^\circ\).由于
\(\angle OAG+\angle GAE=\angle BAE+\angle GAE=90^\circ\)，所以
\(\angle BAE=\angle OAG=40^\circ\)，也即晷针与点 \(A\) 处的水平面所成角为 \(\angle BAE=40^\circ\).

\bitmapfigure[max width=0.62\linewidth,max totalheight=4.8cm]{img/2020/national_paper_2/q4_solution_fig1.png}

故选 B.
\end{solution}

\begin{problem}
某中学的学生积极参加体育锻炼，其中有 \(96\%\) 的学生喜欢足球或游泳，\(60\%\) 的学生喜欢足球，\(82\%\) 的学生喜欢游泳，则该中学既喜欢足球又喜欢游泳的学生数占该校学生总数的比例是（\quad）

\choices
    {62\%}
    {56\%}
    {46\%}
    {42\%}

\end{problem}

\begin{answer}
C
\end{answer}

\begin{solution}
设喜欢足球的学生占比为 \(A\)，喜欢游泳的学生占比为 \(B\)，则
\(P(A)=0.60\)，\(P(B)=0.82\)，\(P(A\cup B)=0.96\).由
\(P(A\cap B)=P(A)+P(B)-P(A\cup B)\)，得
\(P(A\cap B)=0.60+0.82-0.96=0.46\).故选 C.
\end{solution}

\begin{problem}
要安排 \(3\) 名学生到 \(2\) 个乡村做志愿者，每名学生只能选择去一个村，每个村里至少有一名志愿者，则不同的安排方法共有（\quad）

\choices
    {\(2\) 种}
    {\(3\) 种}
    {\(6\) 种}
    {\(8\) 种}

\end{problem}

\begin{answer}
C
\end{answer}

\begin{solution}
第一步，将 \(3\) 名学生分成两个组，有
\(\binom31\binom22=3\)种分法.

第二步，将 \(2\) 组学生安排到 \(2\) 个村，有
\(A_2^2=2\)种安排方法.

所以，不同的安排方法共有
\(3\times2=6\)种.

故选 C.
\end{solution}

\begin{problem}
已知函数 \(f(x)=\lg(x^2-4x-5)\) 在 \((a,+\infty)\) 上单调递增，则 \(a\) 的取值范围是（\quad）

\choices
    {\((2,+\infty)\)}
    {\([2,+\infty)\)}
    {\((5,+\infty)\)}
    {\([5,+\infty)\)}

\end{problem}

\begin{answer}
D
\end{answer}

\begin{solution}
先求定义域：
\(x^2-4x-5>0 \iff x>5 \text{ 或 } x<-1\).
而 \(x^2-4x-5\) 在 \((5,+\infty)\) 上单调递增，所以 \(f(x)=\lg(x^2-4x-5)\) 在 \((5,+\infty)\) 上单调递增.
因此要使 \(f(x)\) 在 \((a,+\infty)\) 上单调递增，必须有
\(a\ge 5\).故选 D.
\end{solution}

\begin{problem}
若定义在 \(\R\) 的奇函数 \(f(x)\) 在 \((-\infty,0)\) 单调递减，且 \(f(2)=0\)，则满足 \(xf(x-1)\ge0\) 的 \(x\) 的取值范围是（\quad）

\choices
    {\([ -1,1]\cup[3,+\infty)\)}
    {\([ -3,-1]\cup[0,1]\)}
    {\([ -1,0]\cup[1,+\infty)\)}
    {\([ -1,0]\cup[1,3]\)}

\end{problem}

\begin{answer}
D
\end{answer}

\begin{solution}
因为 \(f(x)\) 是奇函数，所以
\(f(0)=0\)，\(f(-2)=0\).又因为 \(f(x)\) 在 \((-\infty,0)\) 单调递减，所以

当 \(x\in(-\infty,-2)\) 时，\(f(x)>0\)；

当 \(x\in(-2,0)\) 时，\(f(x)<0\).

由奇函数性质可知 \(f(x)\) 在 \((0,+\infty)\) 也单调递减，于是

当 \(x\in(0,2)\) 时，\(f(x)>0\)；

当 \(x\in(2,+\infty)\) 时，\(f(x)<0\).

因此：

当 \(x\le0\) 时，要使 \(xf(x-1)\ge0\)，需 \(f(x-1)\le0\)，即
\(-1\le x\le0\);
当 \(x\ge0\) 时，要使 \(xf(x-1)\ge0\)，需 \(f(x-1)\ge0\)，即
\(1\le x\le3\).所以所求范围为
\([ -1,0]\cup[1,3]\).故选 D.
\end{solution}
\section{多选题}

\begin{problem}
我国新冠肺炎疫情进入常态化，各地有序推进复工复产，下面是某地连续 \(11\) 天复工复产指数折线图，下列说法正确的是（\quad）

\begin{center}
\bitmapfigure[max width=0.88\linewidth,max totalheight=4.8cm]{img/2020/national_paper_2/q9_fig1.png}
\end{center}

\choices
    {这 \(11\) 天复工指数和复产指数均逐日增加}
    {这 \(11\) 天期间，复产指数增量大于复工指数的增量}
    {第 \(3\) 天至第 \(11\) 天复工复产指数均超过 \(80\%\)}
    {第 \(9\) 天至第 \(11\) 天复产指数增量大于复工指数的增量}

\end{problem}

\begin{answer}
CD
\end{answer}

\begin{solution}
由折线图可知，复工指数在第 \(1\) 天到第 \(2\) 天、第 \(7\) 天到第 \(8\) 天、第 \(10\) 天到第 \(11\) 天都有下降，所以 A 错误.

比较第 \(1\) 天与第 \(11\) 天的数值，复工指数增量大于复产指数增量，所以 B 错误.

第 \(3\) 天至第 \(11\) 天两条折线均在 \(80\%\) 以上，所以 C 正确.

第 \(9\) 天至第 \(11\) 天复产指数的增量大于复工指数的增量，所以 D 正确.
故选 CD.
\end{solution}

\begin{problem}
已知曲线 \(C:mx^2+ny^2=1\). 则下列正确的是（\quad）

\choices
    {若 \(m>n>0\)，则 \(C\) 是椭圆，其焦点在 \(y\) 轴上}
    {若 \(m=n>0\)，则 \(C\) 是圆，其半径为 \(n\)}
    {若 \(mn<0\)，则 \(C\) 是双曲线，其渐近线方程为 \(y=\pm\sqrt{-\dfrac{m}{n}}\,x\)}
    {若 \(m=0\)，\(n>0\)，则 \(C\) 是两条直线}

\end{problem}

\begin{answer}
ACD
\end{answer}

\begin{solution}
对于 A，若 \(m>n>0\)，则
\(mx^2+ny^2=1 \iff \frac{x^2}{1/m}+\frac{y^2}{1/n}=1\).因为 \(m>n\)，所以 \(1/m<1/n\)，故焦点在 \(y\) 轴上，A 正确.

对于 B，若 \(m=n>0\)，则
\(mx^2+ny^2=1 \iff x^2+y^2=\frac1n\)，圆的半径应为 \(\dfrac1{\sqrt n}\)，故 B 错误.

对于 C，若 \(mn<0\)，则 \(C\) 是双曲线，其渐近线为
\(y=\pm\sqrt{-\frac{m}{n}}\,x\)，故 C 正确.

对于 D，若 \(m=0\)，\(n>0\)，则
\(ny^2=1\)，表示两条直线，故 D 正确.
故选 ACD.
\end{solution}

\begin{problem}
下图是函数 \(y=\sin(\omega x+\varphi)\) 的部分图像，则 \(\sin(\omega x+\varphi)=\)（\quad）

\begin{center}
\bitmapfigure[max width=0.42\linewidth,max totalheight=4.8cm]{img/2020/national_paper_2/q11_fig1.png}
\end{center}

\choices
    {\(\sin\left(x+\dfrac{\pi}{3}\right)\)}
    {\(\sin\left(\dfrac{\pi}{3}-2x\right)\)}
    {\(\cos\left(2x+\dfrac{\pi}{6}\right)\)}
    {\(\cos\left(\dfrac{5\pi}{6}-2x\right)\)}

\end{problem}

\begin{answer}
BC
\end{answer}

\begin{solution}
由图象可知周期为
\(T=\pi\)，所以
\(\omega=\frac{2\pi}{T}=2\).又由图象可知在 \(x=\dfrac{\pi}{6}\) 处函数值为 \(0\)，且函数在该点处由正变负，因此
\(2\cdot \frac{\pi}{6}+\varphi=\pi\)，解得
\(\varphi=\frac{2\pi}{3}\).于是
\(\sin(\omega x+\varphi)=\sin\left(2x+\frac{2\pi}{3}\right)\).利用诱导公式可写成
\(\sin\left(2x+\frac{2\pi}{3}\right)=\cos\left(2x+\frac{\pi}{6}\right)=\sin\left(\frac{\pi}{3}-2x\right)\).
故选 BC.
\end{solution}

\begin{problem}
已知 \(a>0\)，\(b>0\)，且 \(a+b=1\)，则（\quad）

\choices
    {\(a^2+b^2\ge\dfrac12\)}
    {\(2^{a-b}>\dfrac12\)}
    {\(\log_2 a+\log_2 b\ge-2\)}
    {\(\sqrt a+\sqrt b\le \sqrt2\)}

\end{problem}

\begin{answer}
ABD
\end{answer}

\begin{solution}
对于 A，
\(a^2+b^2=(a+b)^2-2ab=1-2ab\).因为
\(ab\le \left(\frac{a+b}{2}\right)^2=\frac14\)，所以
\(a^2+b^2\ge 1-\frac12=\frac12\).故 A 正确.

对于 B，因为 \(a-b> -1\)，所以
\(2^{a-b}>2^{-1}=\frac12\).故 B 正确.

对于 C，
\(\log_2 a+\log_2 b=\log_2(ab)\le \log_2\frac14=-2\)，所以 C 错误.

对于 D，
\((\sqrt a+\sqrt b)^2=a+b+2\sqrt{ab}\le a+b+a+b=2\)，故
\(\sqrt a+\sqrt b\le \sqrt2\).故 D 正确.
故选 ABD.
\end{solution}
\section{填空题}

\begin{problem}
已知正方体 \(ABCD-A_1B_1C_1D_1\) 的棱长为 \(2\)，\(M\)、\(N\) 分别为 \(BB_1\)、\(AB\) 的中点，则三棱锥 \(A-NMD_1\) 的体积为\fillinblank{}

\begin{center}
\bitmapfigure[max width=0.52\linewidth,max totalheight=4.8cm]{img/2020/national_paper_2/q13_solution_fig1.png}
\end{center}

\end{problem}

\begin{answer}
\(\dfrac13\)
\end{answer}

\begin{solution}
因为正方体 \(ABCD-A_1B_1C_1D_1\) 的棱长为 \(2\)，\(M\)、\(N\) 分别为 \(BB_1\)、\(AB\) 的中点，
所以
\(V_{A-NMD_1}=V_{D_1-AMN}=\frac13\times\frac12\times1\times1\times2=\frac13\).

故答案为 \(\dfrac13\).
\end{solution}

\begin{problem}
斜率为 \(3\) 的直线过抛物线 \(C:y^2=4x\) 的焦点，且与 \(C\) 交于 \(A\)、\(B\) 两点，则 \(AB=\fillinblank{}\) .

\begin{center}
\bitmapfigure[max width=0.52\linewidth,max totalheight=4.8cm]{img/2020/national_paper_2/q14_solution_fig1.png}
\end{center}

\end{problem}

\begin{answer}
\(\dfrac{16}{3}\)
\end{answer}

\begin{solution}
因为抛物线 \(y^2=4x\) 的焦点为 \(F(1,0)\)，

又因为直线 \(AB\) 过焦点 \(F\) 且斜率为 \(3\)，所以直线 \(AB\) 的方程为
\(y=3(x-1)\).
代入抛物线方程，消去 \(y\) 并化简得
\(3x^2-10x+3=0\).
解法一：解得
\(x_1=\frac13\)，\(x_2=3\).所以
\(|AB|=\sqrt{1+k^2}\,|x_1-x_2|=\sqrt{1+3^2}\left|3-\frac13\right|=\frac{16}{3}\).

解法二：设 \(A(x_1,y_1)\)、\(B(x_2,y_2)\)，则
\(x_1+x_2=\frac{10}{3}\).过 \(A\)、\(B\) 分别作准线 \(x=-1\) 的垂线，设垂足分别为 \(C\)、\(D\). 由抛物线的定义可得
\(|AB|=|AF|+|BF|=|AC|+|BD|=x_1+1+x_2+1=x_1+x_2+2=\frac{16}{3}\).

故答案为 \(\dfrac{16}{3}\).
\end{solution}

\begin{problem}
将数列 \(\{2n-1\}\) 与 \(\{3n-2\}\) 的公共项从小到大排列得到数列 \(\{a_n\}\)，则 \(\{a_n\}\) 的前 \(n\) 项和为\fillinblank{} .

\end{problem}

\begin{answer}
\(3n^2-2n\)
\end{answer}

\begin{solution}
数列 \(\{2n-1\}\) 是首项为 \(1\)、公差为 \(2\) 的等差数列，数列 \(\{3n-2\}\) 是首项为 \(1\)、公差为 \(3\) 的等差数列.
它们的公共项构成首项为 \(1\)、公差为 \(6\) 的等差数列，所以
\(a_n=1+6(n-1)=6n-5\).因此前 \(n\) 项和为
\(S_n=\frac n2\bigl(2\cdot1+(n-1)\cdot6\bigr)=3n^2-2n\).故答案为 \(3n^2-2n\).
\end{solution}

\begin{problem}
某中学开展劳动实习，学生加工制作零件，零件的截面如图所示. \(O\) 为圆孔及轮廓圆弧 \(AB\) 所在圆的圆心，\(A\) 是圆弧 \(AB\) 与直线 \(AG\) 的切点，\(B\) 是圆弧 \(AB\) 与直线 \(BC\) 的切点，四边形 \(DEFG\) 为矩形，\(BC\perp DG\)，垂足为 \(C\)，\(\tan\angle ODC=\dfrac35\)，\(BH\parallel DG\)，\(EF=12\text{ cm}\)，\(DE=2\text{ cm}\)，\(A\) 到直线 \(DE\) 和 \(EF\) 的距离均为 \(7\text{ cm}\)，圆孔半径为 \(1\text{ cm}\)，则图中阴影部分的面积为\fillinblank{} \(\text{cm}^2\).

\begin{center}
\bitmapfigure[max width=0.88\linewidth,max totalheight=4.8cm]{img/2020/national_paper_2/q16_fig1.png}
\end{center}

\end{problem}

\begin{answer}
\(4+\dfrac{5\pi}{2}\)
\end{answer}

\begin{solution}
设 \(OB=OA=r\)，由题意 \(AM=AN=7\)，\(EF=12\)，所以 \(NF=5\)；

\begin{center}
\bitmapfigure[max width=0.60\linewidth,max totalheight=4.8cm]{img/2020/national_paper_2/q20_solution_fig1.png}
\end{center}

因为 \(AP=5\)，所以 \(\angle AGP=45^\circ\)；

因为 \(BH//DG\)，所以 \(\angle AHO=45^\circ\)；

因为 \(AG\) 与圆弧 \(AB\) 相切于 \(A\) 点，所以 \(OA\perp AG\)，即 \(\triangle OAH\) 为等腰直角三角形.

在直角 \(\triangle OQD\) 中，
\(OQ=5-\frac{\sqrt2}{2}r\)，\(DQ=7-\frac{\sqrt2}{2}r\).因为 \(\tan\angle ODC=\dfrac35\)，所以
\(\frac{OQ}{DQ}=\frac35\)，即
\(21-3\sqrt2\,r=25-5\sqrt2\,r\)，解得
\(r=2\sqrt2\).
等腰直角 \(\triangle OAH\) 的面积为
\(S_1=\frac12\times2\sqrt2\times2\sqrt2=4\).
扇形 \(AOB\) 的面积
\(S_2=\frac12\times\frac{3\pi}{4}\times\left( 2\sqrt2 \right)^2=3\pi\).
所以阴影部分的面积为
\(S_1+S_2-\frac12\pi=4+\frac{5\pi}{2}\).
故答案为 \(4+\dfrac{5\pi}{2}\).
\end{solution}
\section{解答题}

\begin{problem}
在
\begin{circlelist}
\item \(ac=\sqrt3\)
\item \(c\sin A=3\)
\item \(c=\sqrt3\,b\)
\end{circlelist}

这三个条件中任选一个，补充在下面问题中，若问题中的三角形存在，求 \(c\) 的值；若问题中的三角形不存在，说明理由.

问题：是否存在 \(\triangle ABC\)，它的内角 \(A,B,C\) 的对边分别为 \(a,b,c\)，且 \(\sin A=\sqrt3\sin B\)，\(C=\dfrac{\pi}{6}\)，\fillinblank{}？
\end{problem}

\begin{solution}
由正弦定理得
\(\frac a b=\frac{\sin A}{\sin B}=\sqrt3\).不妨设
\(a=\sqrt3\,m\)，\(b=m\)（\(m>0\)）.又因为 \(C=\dfrac{\pi}{6}\)，由余弦定理得
\(c^2=a^2+b^2-2ab\cos C=3m^2+m^2-2\cdot\sqrt3 m\cdot m\cdot \frac{\sqrt3}{2}=m^2\)，
即
\(c=m\).
\text{解法一：}

（1）若选（1） \(ac=\sqrt3\)，则
\(\sqrt3\,m\cdot m=\sqrt3\)，解得 \(m=1\)，所以
\(c=1\).
（2）若选（2） \(c\sin A=3\)，则
\(\sin A=\sin\frac{2\pi}{3}=\frac{\sqrt3}{2}\)，故
\(c\cdot \frac{\sqrt3}{2}=3\)，解得
\(c=2\sqrt3\).
（3）若选（3） \(c=\sqrt3\,b\)，则
\(\frac cb=\frac mm=1\)，即 \(c=b\)，与条件 \(c=\sqrt3\,b\) 矛盾，所以问题中的三角形不存在.

\text{解法二：} 因为 \(\sin A=\sqrt3\sin B\)，\(C=\dfrac{\pi}{6}\)，\(B=\pi-\left( A+C \right)\)，

所以
\(\sin A=\sqrt3\sin\left( A+C \right)=\sqrt3\sin\left( A+\frac{\pi}{6} \right)\).
于是
\(\sin A=\sqrt3\sin\left( A+C \right)=\sqrt3\sin A\cdot\frac{\sqrt3}{2}+\sqrt3\cos A\cdot\frac12\)，即
\(\sin A=-\sqrt3\cos A\).所以
\(\tan A=-\sqrt3\)，从而
\(A=\frac{2\pi}{3}\)，\(B=C=\frac{\pi}{6}\).
若选（1），因为 \(ac=\sqrt3\)，\(a=\sqrt3\,b=\sqrt3\,c\)，所以
\(\sqrt3\,c^2=\sqrt3\)，故
\(c=1\).
若选（2），因为 \(c\sin A=3\)，所以
\(\frac{\sqrt3 c}{2}=3\)，故
\(c=2\sqrt3\).
若选（3），则与条件 \(c=\sqrt3\,b\) 矛盾.
\end{solution}

\begin{problem}
已知公比大于 \(1\) 的等比数列 \(\{a_n\}\) 满足 \(a_2+a_4=20\)，\(a_3=8\).

\begin{enumerate}
\item 求 \(\{a_n\}\) 的通项公式；
\item 求 \(a_1a_2-a_2a_3+\cdots+(-1)^{n-1}a_na_{n+1}\).
\end{enumerate}

\end{problem}

\begin{solution}
\text{（1）} 设等比数列 \(\{a_n\}\) 的公比为 \(q(q>1)\)，首项为 \(a_1\)，则
\examdisplaycases{}{
a_1q+a_1q^3=20,\\
a_1q^2=8.
}
由第二式得 \(a_1=\dfrac8{q^2}\)，代入第一式得
\(\frac8{q^2}(q+q^3)=20\)，即
\(2q^2-5q+2=0\).解得
\(q=2\)（\(\text{舍去 }q=\tfrac12\)），从而
\(a_1=2\).\(a_n=2\cdot2^{n-1}=2^n\).
\text{（2）} 由于
\((-1)^{n-1}a_na_{n+1}=(-1)^{n-1}\cdot2^n\cdot2^{n+1}=(-1)^{n-1}2^{2n+1}\)，故
\(a_1a_2-a_2a_3+\cdots+(-1)^{n-1}a_na_{n+1}=2^3-2^5+2^7-2^9+\cdots+(-1)^{n-1}2^{2n+1}\).
所以
\(S_n=\frac{2^3\bigl( 1-(-4)^n \bigr)}{1-(-4)}=\frac85-\frac{2^{2n+3}}{5}(-1)^n\).
\end{solution}

\begin{problem}
为加强环境保护，治理空气污染，环境监测部门对某市空气质量进行调研，随机抽查了 \(100\) 天空气中的 PM\(_{2.5}\) 和 SO\(_2\) 浓度（单位：\(\mu\text{g}/\text{m}^3\)），得下表：

\begin{center}
\begin{tabular}{c|ccc}
\hline
PM\(_{2.5}\)\(\backslash\)SO\(_2\) & \([0,50]\) & \((50,150]\) & \((150,475]\) \\
\hline
\([0,35]\) & 32 & 18 & 4 \\
\((35,75]\) & 6 & 8 & 12 \\
\((75,115]\) & 3 & 7 & 10 \\
\hline
\end{tabular}
\end{center}

\begin{enumerate}
\item 估计事件“该市一天空气中 PM\(_{2.5}\) 浓度不超过 \(75\)，且 SO\(_2\) 浓度不超过 \(150\)”的概率；
\item 根据所给数据，完成下面的 \(2\times2\) 列联表；
\item 根据（2）中的列联表，判断是否有 \(99\%\) 的把握认为该市一天空气中 PM\(_{2.5}\) 浓度与 SO\(_2\) 浓度有关？
\end{enumerate}

\end{problem}

\begin{solution}
（1）满足条件的天数为
\(32+6+18+8=64\)，所以所求概率为
\(\frac{64}{100}=0.64\).
（2）由题中数据可得列联表如下：
\examdisplayarray{}{c|cc|c}{
\hline
& [0,150] & (150,475] & \text{合计} \\
\hline
[0,75] & 64 & 16 & 80 \\
(75,115] & 10 & 10 & 20 \\
\hline
\text{合计} & 74 & 26 & 100 \\
\hline
}{}

（3）由列联表中的数据可得
\(K^2=\frac{100(64\cdot10-16\cdot10)^2}{80\cdot20\cdot74\cdot26}=\frac{3600}{481}\approx7.4844\).
因为 \(7.484>6.635\)，所以有 \(99\%\) 的把握认为该市一天空气中 PM\(_{2.5}\) 浓度与 SO\(_2\) 浓度有关.
\end{solution}

\begin{problem}
如图，四棱锥 \(P-ABCD\) 底面为正方形，\(PD\perp\) 底面 \(ABCD\). 设平面 \(PAD\) 与平面 \(PBC\) 的交线为 \(l\).



\begin{enumerate}
\item 证明：\(l\perp\) 平面 \(PDC\)；
\item 已知 \(PD=AD=1\)，\(Q\) 为 \(l\) 上的点，\(QB=\sqrt2\)，求 \(PB\) 与平面 \(QCD\) 所成角的正弦值.
\end{enumerate}

\bitmapfigure[max width=0.88\linewidth,max totalheight=4.8cm]{img/2020/national_paper_2/q20_fig1.png}

\end{problem}

\begin{solution}
（1）在正方形 \(ABCD\) 中，
\(AD\parallel BC\).又因为 \(BC\subset\) 平面 \(PBC\)，且 \(AD\not\subset\) 平面 \(PBC\)，所以
\(AD\parallel \text{平面 }PBC\).而 \(AD\subset\) 平面 \(PAD\)，且
\(\text{平面 }PAD\cap\text{平面 }PBC=l\)，所以
\(AD\parallel l\).又因为正方形 \(ABCD\) 中 \(AD\perp DC\)，且 \(PD\perp\) 底面 \(ABCD\)，所以
\(AD\perp DC\)，\(AD\perp PD\).故 \(AD\perp\) 平面 \(PDC\)，从而 \(l\perp\) 平面 \(PDC\).

（2）建立空间直角坐标系 \(D-xyz\)，设
\(D(0,0,0)\)，\(C(0,1,0)\)，\(A(1,0,0)\)，\(P(0,0,1)\)，\(B(1,1,0)\).设 \(Q(m,0,1)\)，则
\(\overrightarrow{DC}=(0,1,0)\)，\(\overrightarrow{DQ}=(m,0,1)\)，\(\overrightarrow{PB}=(1,1,-1)\).由 \(QB=\sqrt2\) 得
\((m-1)^2+(0-1)^2+(1-0)^2=2\)，解得
\(m=1\).设平面 \(QCD\) 的法向量为 \(\bs n=(x,y,z)\)，则
\(\overrightarrow{DC}\cdot\bs n=0\)，\(\overrightarrow{DQ}\cdot\bs n=0\).于是
\(y=0\)，\(x+z=0\).令 \(x=1\)，则 \(z=-1\)，所以可取
\(\bs n=(1,0,-1)\).于是
\(\cos\angle(\bs n,\overrightarrow{PB})=\frac{\left|\bs n\cdot\overrightarrow{PB}\right|}{|\bs n|\cdot|\overrightarrow{PB}|}=\frac{|1+1|}{\sqrt2\cdot\sqrt3}=\frac{\sqrt6}{3}\).
而直线与平面所成角的正弦值等于其方向向量与平面法向量所成角的余弦值的绝对值，所以
\(\sin\angle\bigl(PB,\text{平面 }QCD\bigr)=\frac{\sqrt6}{3}\).
\end{solution}

\begin{problem}
已知椭圆 \(C:\dfrac{x^2}{a^2}+\dfrac{y^2}{b^2}=1\)（\(a>b>0\)） 过点 \(M(2,3)\)，点 \(A\) 为其左顶点，且 \(AM\) 的斜率为 \(\dfrac12\)，



\begin{enumerate}
\item 求 \(C\) 的方程；
\item 点 \(N\) 为椭圆上任意一点，求 \(\triangle AMN\) 的面积的最大值.
\end{enumerate}

\bitmapfigure[max width=0.52\linewidth,max totalheight=4.8cm]{img/2020/national_paper_2/q21_solution_fig1.png}

\end{problem}

\begin{solution}
（1）因为 \(A\) 为椭圆左顶点，所以 \(A(-a,0)\). 由 \(AM\) 的斜率为 \(\dfrac12\)，得
\(\frac{3-0}{2-(-a)}=\frac12\)，解得
\(a=4\).又因为点 \(M(2,3)\) 在椭圆上，所以
\(\frac{4}{16}+\frac{9}{b^2}=1\)，解得
\(b^2=12\).故椭圆方程为
\(\frac{x^2}{16}+\frac{y^2}{12}=1\).
（2）直线 \(AM\) 的方程为
\(y-3=\frac12(x-2)\)，即
\(x-2y+4=0\).设与直线 \(AM\) 平行的直线方程为
\(x-2y+m=0\).当该直线与椭圆相切时，点 \(N\) 到直线 \(AM\) 的距离最大，从而 \(\triangle AMN\) 的面积最大.

将 \(x=2y-m\) 代入椭圆方程得
\(\frac{(2y-m)^2}{16}+\frac{y^2}{12}=1\).整理得
\(16y^2-12my+3m^2-48=0\).切线条件为判别式为零，即
\((-12m)^2-4\cdot16(3m^2-48)=0\)，解得
\(m=\pm8\).与 \(AM\) 距离更远的平行为 \(x-2y+8=0\). 故点 \(N\) 到直线 \(AM\) 的最大距离为两平行线间距离
\(d=\frac{|8-(-4)|}{\sqrt{1^2+(-2)^2}}=\frac{12}{\sqrt5}\).又
\(|AM|=\sqrt{(2+4)^2+3^2}=3\sqrt5\).所以 \(\triangle AMN\) 的面积最大值为
\(\frac12\cdot3\sqrt5\cdot\frac{12}{\sqrt5}=18\).
\end{solution}

\begin{problem}
已知函数 \(f(x)=ae^{x-1}-\ln x+\ln a\).

\begin{enumerate}
\item 当 \(a=e\) 时，求曲线 \(y=f(x)\) 在点 \((1,f(1))\) 处的切线与两坐标轴围成的三角形的面积；
\item 若 \(f(x)\ge1\)，求 \(a\) 的取值范围.
\end{enumerate}

\end{problem}

\begin{solution}
（1）当 \(a=e\) 时，
\(f(x)=e^x-\ln x+1\).于是
\(f'(x)=e^x-\frac1x\).所以
\(f(1)=e+1\)，\(f'(1)=e-1\).切线方程为
\(y-(e+1)=(e-1)(x-1)\)，即
\(y=(e-1)x+2\).与两坐标轴的交点分别为 \((0,2)\) 和 \(\left(-\dfrac2{e-1},0\right)\)，故三角形面积为
\(\frac12\cdot2\cdot\frac2{e-1}=\frac{2}{e-1}\).
（2）因为
\(f(x)=ae^{x-1}-\ln x+\ln a\)，所以
\(f'(x)=ae^{x-1}-\frac1x\)，且 \(a>0\).

设 \(g(x)=f'(x)\)，则
\(g'(x)=ae^{x-1}+\frac1{x^2}>0\).所以 \(g(x)\) 在 \((0,+\infty)\) 上单调递增，即 \(f'(x)\) 在 \((0,+\infty)\) 上单调递增.

当 \(a=1\) 时，
\(f'(1)=0\)，所以
\(f(x)_{\min}=f(1)=1\)，故 \(f(x)\ge1\) 成立.

当 \(a>1\) 时，因为
\(\frac1a<1\)，所以
\(e^{\frac1a-1}<1\).于是
\(f'\left( \frac1a \right)f'(1)=a\left( e^{\frac1a-1}-1 \right)(a-1)<0\).
故存在唯一 \(x_0>0\)，使得
\(f'(x_0)=ae^{x_0-1}-\frac1{x_0}=0\)，且当 \(x\in\left( 0,x_0 \right)\) 时，\(f'(x)<0\)；当 \(x\in\left( x_0,+\infty \right)\) 时，\(f'(x)>0\).

于是
\(ae^{x_0-1}=\frac1{x_0}\)，从而
\(\ln a+x_0-1=-\ln x_0\).因此 \(f(x)_{\min}=f(x_0)=ae^{x_0-1}-\ln x_0+\ln a=\frac1{x_0}+\ln a+x_0-1+\ln a\ge2\ln a-1+2\sqrt{\frac1{x_0}\cdot x_0}=2\ln a+1>1\).
故 \(f(x)\ge1\) 恒成立.

当 \(0<a<1\) 时，
\(f(1)=a+\ln a<a<1\)，故 \(f(x)\ge1\) 不是恒成立.

综上，实数 \(a\) 的取值范围是
\([1,+\infty)\).
\end{solution}
